\chapter[Average Hessian Operator]{The Average Hessian Operator and Universal Bregman Theory}\label{ch:sec:the-average-hessian-operator-and-universal-bregman-theo}

The Average Hessian Operator of this chapter turns the universal inequalities of Chapter~15 into exact equalities for Bregman divergences. This operator plays a central role in the variational theory of Part~V, where it connects the UOD to the spectral geometry of the posterior Laplacian.


\section{Introduction}\label{ch:sec:introduction-bregman}
The universal Lipschitz and quadratic bounds apply to any smooth
functional.  But some functionals---Bregman divergences---have a special
structure: they admit an exact quadratic representation through an
average of their Hessian along the geodesic.  For these functionals the
bounds become exact identities, not estimates.  This chapter develops
the Average Hessian Operator, which turns the universal inequalities of
the previous chapter into precise equalities for the entire Bregman
class.

Chapter~\ref{ch:sec:information-functionals-on-the-wis-manifold} established geometric estimates that hold for arbitrary
smooth scalar fields on the WIS manifold. We now specialize those
results to the important class of \textbf{Bregman divergences}.
The central idea of this chapter is not merely that Bregman divergences
admit an integral representation---a classical result---but that, within
the WIS framework, this representation naturally defines a canonical
symmetric operator associated with every pair of posterior
distributions.
This operator, called the \textbf{Average Hessian Operator}, provides the
bridge between the universal geometry induced by the Universal
Optimality Defect (UOD) and the anisotropic behavior of individual
information-theoretic divergences.
\subsection{Convex Potentials on the WIS Coordinate Space}\label{ch:sec:convex-potentials-on-the-wis-coordinate-space}
Let
$F:\mathcal P\rightarrow F(\mathcal P)$
be the global WIS coordinate map established in Part~\ref{ch:part:posterior}.
Let $\phi:F(\mathcal P)\rightarrow\mathbb R$ be a twice continuously differentiable convex potential.
For
$x=F(P),\qquad y=F(Q),$
the associated Bregman divergence is
\[ D_\phi(P,Q) =
\phi(x)-\phi(y)-\nabla\phi(y)^T(x-y).
\]
Throughout this chapter all line segments are understood in the WIS
coordinate space \(F(\mathcal P)\).

\begin{theorem}[Integral Representation]\label{ch:thm:15-1}

If \(\phi\in C^{2}(F(\mathcal P))\), then
\[ D_\phi(P,Q) = \int_0^{1} (1-t)\,\Delta^T
H_\phi(\gamma(t)) \Delta\,dt, \]
where
$H_\phi=\nabla^2\phi.$
\end{theorem}

\paragraph{Interpretation.}
The Average Hessian Operator turns an inequality into an identity for Bregman divergences.  While the universal bounds of Chapter~15 give estimates of the form $|S(P)-S(Q)|\le L\sqrt{\mathcal D}$, the AHO gives exact equalities $D_\phi(P,Q)=\frac12\Delta^T\overline H_\phi\Delta$.  This is the bridge between the universal theory and the specific geometry of each functional.

\begin{quote}
\textbf{Mathematical intuition.}
The Average Hessian Operator $\overline H_\phi(P,Q)$ averages the Hessian of $\phi$ along the geodesic joining $P$ to $Q$.
For Bregman divergences this average turns an inequality into an identity: $D_\phi(P,Q)=\frac12\Delta^T\overline H_\phi\Delta$.
The UOD bounds the spectral norm of this operator, giving universal estimates that depend only on the geometry.
\end{quote}

\begin{proof}
Taylor's theorem with integral remainder gives
\[
\phi(x) - \phi(y) - \nabla\phi(y)^T(x-y)
= \int_0^{1} (1-t)\,(x-y)^T H_\phi(y+t(x-y))\,(x-y)\,dt.
\]
The left-hand side is the Bregman divergence $D_\phi(P,Q)$ by definition,
and the right-hand side is precisely the integral expression claimed. 
\end{proof}



\section{The Average Hessian Operator}\label{ch:sec:the-average-hessian-operator}
The previous theorem motivates the following definition.
\begin{definition}[Average Hessian Operator]
\label{ch:defn:15-1}
The \textbf{Average Hessian Operator} associated with \(\phi\) between
\(P\) and \(Q\) is
\[ \boxed{
\overline H_\phi(P,Q)
=
2\int_0^1
(1-t)
H_\phi(\gamma(t))\,dt.
} \]
Because every Hessian is symmetric,
$\overline H_\phi(P,Q)$
is symmetric.
Since every Hessian of a convex potential is positive semidefinite,
$\overline H_\phi(P,Q)\succeq0.$
\end{definition}

\paragraph{Why this operator.}
The Average Hessian Operator turns the universal inequality into an exact equality for Bregman divergences.  It is the bridge between the universal theory and the specific geometry of each divergence.

The integral representation immediately becomes
\begin{theorem}[Exact Quadratic Representation]
\label{ch:thm:15-2}
\[ \boxed{
D_\phi(P,Q)
=
\frac12
\Delta^T
\overline H_\phi(P,Q)
\Delta.
} \]
\end{theorem}
Thus every Bregman divergence is represented exactly by a quadratic form
whose metric tensor is the Average Hessian Operator.
\subsection{Spectral Representation}\label{ch:cor:15-1}
Since
$\overline H_\phi(P,Q)$
is real symmetric, the spectral theorem implies
$\overline H_\phi = U\Lambda U^T,$
where \(U\) is orthogonal and
\[
\Lambda=\operatorname{diag}(\lambda_1,\ldots,\lambda_n).
\]
Writing
$\alpha=U^T\Delta,$
gives the exact decomposition
\begin{corollary}[Spectral Decomposition]

\[ \boxed{
D_\phi(P,Q)
=
\frac12
\sum_{i=1}^{n}
\lambda_i\alpha_i^2.
} \]
\end{corollary}
This expression separates the universal displacement from the
directional amplification induced by the convex potential.
\begin{theorem}[Spectral Bounds]\label{ch:thm:15-3}
Assume that along the segment
$\gamma([0,1]),$
the Hessian satisfies
\[ m(t)I \preceq H_\phi(\gamma(t))
\preceq M(t)I. \]
Define
$\underline\lambda = 2\int_0^{1}(1-t)m(t)\,dt,$
and
$\overline\lambda = 2\int_0^{1}(1-t)M(t)\,dt.$
Then
\[ \boxed{
\frac{\underline\lambda}{2}\,\|\Delta\|^2
\;\le\; D_\phi(P,Q)
\;\le\; \frac{\overline\lambda}{2}\,\|\Delta\|^2.
} \]
\end{theorem}
\begin{proof}
The matrix inequalities imply identical inequalities for every quadratic
form. Integrating them over the segment and using the definition of
\(\phi\) yields the result. 
\end{proof}


\section{Lower Bounds for General Smooth Functionals}\label{ch:sec:lower-bounds-for-general-smooth-functionals}

\paragraph{Running example: differential privacy.}
The KL divergence $D_{\mathrm{KL}}(P\|Q)=\sum_i P_i\log(P_i/Q_i)$ between posterior and uniform prior is a Bregman divergence generated by the negative Shannon entropy $\phi(x)=\sum_i x_i\log x_i$.  In the differential privacy example (Section~\ref{ch:sec:differential-privacy-and-database-queries}), the Average Hessian Operator $\overline H_\phi(P,Q)$ is diagonal with eigenvalues $\lambda_i=1/(\pi_i+P_i)$, giving the spectral quadratic form $\Delta^T\overline H_\phi\Delta=\sum_i\Delta_i^2/(\pi_i+P_i)$.  This yields explicit spectral bounds on how much the privacy---utility trade-off can differ between two nearby mechanisms.

The spectral bounds of Theorem~\ref{ch:thm:15-3} apply only to Bregman
divergences, whose quadratic structure is forced by the convex
potential.  For arbitrary smooth functionals $S:\mathcal P\to\mathbb R$,
the situation is different: the second-order Taylor expansion gives
only an upper bound in general (Theorem~\ref{ch:thm:14-2}), and a
matching lower bound requires additional structure.

\subsection{Local Lower Bounds via Positive Definite Hessian}\label{ch:sec:local-lower-bounds-via-positive-definite-hessian}

If $S$ has a stationary point $P^\star$ and its WIS Hessian is positive
definite in a neighbourhood, then $S$ is locally strictly convex and

\[
S(P) - S(P^\star) \ge \frac{\lambda_{\min}}{2}\,\mathcal D(P,P^\star),
\]

where $\lambda_{\min}>0$ is the minimum eigenvalue of the Hessian on the
segment joining $P$ and $P^\star$.

\begin{theorem}[Local Lower Bound]
\label{ch:thm:lower-local}
Let $S\in C^2(\mathcal P)$ have a stationary point $\nabla_g S(P^\star)=0$.
Assume that the WIS Hessian satisfies
$\operatorname{Hess}_g S(x) \succeq \mu I$ for every $x$ on the
geodesic segment joining $P$ and $P^\star$, with $\mu>0$.  Then

\[
S(P) - S(P^\star) \ge \frac{\mu}{2}\,\mathcal D(P,P^\star).
\]

Equivalently, for any $P,Q$ in a convex neighbourhood where
$\operatorname{Hess}_g S \succeq \mu I$ along the segment,

\[
|S(P)-S(Q)| \ge \frac{\mu}{2}\,\mathcal D(P,Q)
\quad\text{(up to sign depending on convexity/concavity)}.
\]

\end{theorem}

\begin{proof}
Taylor's theorem with integral remainder gives

\[
S(P) - S(P^\star)
=
\int_0^1 (1-t)\,\Delta^T\operatorname{Hess}_g S(\gamma(t))\,\Delta\,dt,
\]

where $\gamma$ is the geodesic from $P^\star$ to $P$ and
$\Delta$ is the coordinate difference in the WIS embedding.
Since $\operatorname{Hess}_g S \succeq \mu I$ along the segment,
each quadratic form is bounded below by $\mu\|\Delta\|^2$.  Integrating
and using $\mathcal D(P,P^\star)=\|\Delta\|^2$ gives the result.\end{proof}

\subsection{Global Lower Bounds via Geodesic Convexity}\label{ch:sec:global-lower-bounds-via-geodesic-convexity}

When $S$ is geodesically convex on the entire posterior manifold, the
local lower bound extends globally.  Geodesic convexity---convexity
along geodesics of the pullback metric---is a stronger condition than
Euclidean convexity in coordinates, but it is satisfied by several
important functionals.

\begin{lemma}[Geodesic Convexity of the KL Divergence]
\label{ch:lem:kl-convex}
The Kullback--Leibler divergence $P\mapsto D_{\mathrm{KL}}(P\|Q)$ is
geodesically convex on the posterior manifold for every fixed $Q$.
\end{lemma}

\begin{proof}
In the WIS coordinate representation,
$D_{\mathrm{KL}}(P\|Q) = \sum_i P_i\log(P_i/Q_i)$ becomes
$D_{\mathrm{KL}}(x\|y) = \sum_i e^{x_i-n}(x_i-y_i)$ after the
logarithmic embedding and normalisation.  The Hessian of this
expression in the WIS coordinates is
$\operatorname{diag}\bigl(e^{x_i-n}(x_i-y_i+2)\bigr)$,
which is positive definite precisely when $x_i-y_i+2>0$ for every $i$, in
particular on a neighbourhood of $Q$ (where $x_i=y_i$).  Since the pullback
metric is flat, this implies geodesic convexity on that neighbourhood.
\end{proof}

\paragraph{What this theorem proves.}
The universal bounds of Chapter~15 give upper estimates for how much a functional can change.  Lower bounds are harder because they require convexity.  This theorem shows that for convex functionals, the quadratic lower bound is exactly $\frac12\lambda_{\min}\|\Delta\|^2$, where $\lambda_{\min}$ is the minimum eigenvalue of the Hessian.  This completes the two-sided estimate: both upper and lower bounds now have explicit geometric constants.

\begin{theorem}[Global Lower Bound for Convex Functionals]
\label{ch:thm:lower-global}
Let $S\in C^2(\mathcal P)$ be geodesically convex and have a stationary
point $P^\star$ (necessarily the global minimiser).  Then for every $P$,

\[
S(P) - S(P^\star) \ge \frac{\mu}{2}\,\mathcal D(P,P^\star),
\]

where $\mu$ is the minimum eigenvalue of the Hessian of $S$ over the
entire segment from $P^\star$ to $P$.
\end{theorem}

\subsection{What the Lower Bound Implies}\label{ch:sec:what-the-lower-bound-implies}

The existence of a lower bound for convex functionals complements the
universal upper bounds of Chapter~\ref{ch:sec:information-functionals-on-the-wis-manifold} and the Bregman theory of this
chapter.  Together they bracket every geodesically convex functional
between two multiples of the UOD:

\[
\frac{\mu}{2}\,\mathcal D(P,P^\star)
\le
S(P) - S(P^\star)
\le
\frac{M}{2}\,\mathcal D(P,P^\star),
\]

where $M$ is the maximum eigenvalue of the Hessian along the segment
(Theorem~\ref{ch:thm:14-2}).  For Bregman divergences, $\mu$ and $M$
are the extremal eigenvalues of the Average Hessian Operator
(Theorem~\ref{ch:thm:15-3}).

For non-convex functionals---such as the R\'enyi entropy for
$\alpha>1$, which is concave---the lower bound is reversed (the
functional is bounded above by the quadratic estimate and below by the
trivial constant).  In such cases a lower bound in terms of the UOD is
impossible without additional structure, because the functional can be
arbitrarily flat near its stationary point.

\subsection{Lower Bound for R\'enyi Entropy ($\alpha<1$)}\label{ch:sec:lower-bound-for-renyi-entropy}

For $\alpha<1$, the R\'enyi entropy
$H_\alpha(P)=\frac1{1-\alpha}\log\sum_i P_i^\alpha$ is strictly
concave on the interior simplex.  Its WIS Hessian, computed in
Remark~\ref{rem:improved-hessian}, is negative definite:

\[
\operatorname{Hess}_g H_\alpha(P) = \frac{\alpha}{1-\alpha}
\frac{\operatorname{diag}(P^{\alpha-2})}{\sum_i P_i^\alpha}
- \frac{\alpha^2}{(1-\alpha)^2}
\frac{P^{\alpha-1}(P^{\alpha-1})^T}{(\sum_i P_i^\alpha)^2}
\prec 0.
\]

The uniform distribution $U_n$ is the unique maximiser (the
stationary point of a concave functional).  Applying the lower bound
argument with reversed sign gives

\[
H_\alpha(U_n) - H_\alpha(P) \ge \frac{|\mu_\alpha|}{2}\,
\mathcal D(P,U_n),
\]

where $|\mu_\alpha|>0$ is the minimum absolute eigenvalue of the
Hessian on the segment.  For $P$ close to $U_n$, this simplifies to

\[
\log_2 n - H_\alpha(P) \ge \frac{\alpha}{2(1-\alpha)}\,
\frac{1}{n}\,\mathcal D(P,U_n) + o\bigl(\mathcal D(P,U_n)\bigr).
\]

Thus the R\'enyi entropy for $\alpha<1$ is bounded below by the
same quadratic form that bounds it above for $\alpha>1$, with the
sign reversed.

\subsection{Lower Bound for Min-Entropy}\label{ch:sec:lower-bound-for-min-entropy}

Min-entropy $H_\infty(P)=-\log\max_i P_i$ is concave on each
regularity region (where the argmax is fixed).  On such a region,
its WIS Hessian is negative semidefinite with a unique negative
eigenvalue.  Consequently,

\[
\log_2 n - H_\infty(P) \ge \frac{1}{2(\max_i P_i)^2}\,
\mathcal D(P,U_n)
\]

for distributions sufficiently close to the uniform distribution
(where the argmax is unique and constant).  This lower bound is the
counterpart of the upper bound derived in
Corollary~\ref{ch:cor:g-entropy-lipschitz} and Theorem~\ref{ch:thm:max-leakage}.

\subsection{Lower Bound for Jensen--Shannon Divergence}\label{ch:sec:lower-bound-for-jensen--shannon-divergence}

The Jensen--Shannon divergence
$\operatorname{JSD}(P,Q)$ is not a Bregman divergence, but it can be
expressed as the average of two KL divergences:

\[
\operatorname{JSD}(P,Q)
= \frac12 D_{\mathrm{KL}}(P\|M) + \frac12 D_{\mathrm{KL}}(Q\|M),
\qquad M=\frac{P+Q}{2}.
\]

Each KL divergence is geodesically convex (Lemma~\ref{ch:lem:kl-convex}),
and the average preserves convexity.  Therefore JSD is geodesically
convex in each argument separately, and for a fixed reference $Q$,

\[
\operatorname{JSD}(P,Q) \ge \frac{\mu_{\mathrm{JSD}}}{2}\,
\mathcal D(P,Q),
\]

where $\mu_{\mathrm{JSD}}$ is the minimum eigenvalue of the
average of the two KL Hessians along the segment.  For $P$ and $Q$
close to the uniform distribution $U_n$, the leading-order
approximation gives

\[
\operatorname{JSD}(P,Q) \ge \frac{1}{4n}\,\mathcal D(P,Q) + o\bigl(\mathcal D(P,Q)\bigr).
\]

\subsection{Unified Bracketing}\label{ch:sec:unified-bracketing}

The lower bounds derived here, together with the upper bounds of
Chapter~\ref{ch:sec:information-functionals-on-the-wis-manifold} and the spectral bounds of Theorem~\ref{ch:thm:15-3},
give a complete bracketing for all major information functionals:

\[
\begin{array}{lll}
\textbf{Functional} & \textbf{Lower bound} & \textbf{Upper bound} \\ \hline
\text{KL divergence} & \frac{\mu}{2}\,\mathcal D & \frac{M}{2}\,\mathcal D \\
\text{R\'enyi entropy }\alpha>1 & \text{none (concave)} & C_\alpha\,\mathcal D \\
\text{R\'enyi entropy }\alpha<1 & \frac{|\mu_\alpha|}{2}\,\mathcal D & \text{none (convex)} \\
\text{Min-entropy} & \frac12(\max_i P_i)^{-2}\,\mathcal D & \frac12(\min_i P_i)^{-2}\,\mathcal D \\
\text{Jensen--Shannon} & \frac{1}{4n}\,\mathcal D & \frac12\,\mathcal D \\
\text{Bregman divergences} & \frac{\lambda_{\min}}{2}\,\mathcal D & \frac{\lambda_{\max}}{2}\,\mathcal D
\end{array}
\]

where the constants are evaluated on the relevant compact neighbourhoods
and $\mu,M$ denote the extremal Hessian eigenvalues along the segment.

\begin{theorem}[Comparison Principle]\label{ch:thm:15-4}

Let \(\phi\) and \(\psi\) be convex potentials satisfying
$H_\phi(x)\preceq H_\psi(x)$
for every point of the WIS segment joining \(P\) and \(Q\).
Then
\[ \boxed{
D_\phi(P,Q)\le D_\psi(P,Q).
} \]
\end{theorem}
\begin{proof}
Subtract the integral representations of the two divergences. Since the
difference of Hessians is positive semidefinite at every point of the
segment, the integrand is non-negative, and so is its integral.
\end{proof}


\section{Relation with the Universal Optimality Defect}\label{ch:sec:relation-with-the-universal-optimality-defect}
The Universal Optimality Defect is
$\mathcal D(P,Q)=\|\Delta\|^2.$
It is therefore the \textbf{canonical isotropic quadratic reference form}
associated with the WIS geometry.
Every Bregman divergence is obtained from the same displacement vector
by replacing the identity operator with the Average Hessian Operator,
\[ D_\phi(P,Q) = \frac12 \Delta^T
\overline H_\phi(P,Q) \Delta. \]
Consequently, the geometry of the WIS manifold is universal, while
individual divergences differ only through the operator
\(\overline H_\phi(P,Q)\).
As a special case,
$\overline H_\phi=I$
implies
$D_\phi(P,Q) = \frac12 \mathcal D(P,Q).$
\paragraph{Outlook.}
The Average Hessian Operator provides a common mathematical object
through which diverse information-theoretic divergences can be analyzed.
The following chapter applies this framework to classical smooth
information functionals, showing that Shannon entropy, Kullback--Leibler
divergence, R\'enyi entropy, Tsallis entropy, collision entropy, and
related quantities become instances of the same geometric theory.
